\chapter{Database Relationships}

\section{One-to-One: User and Profile}

\begin{diagrambox}[One-to-One]
\begin{verbatim}
User Document                    Profile Document
--------------                   -----------------
_id: "u1"                        _id: "p1"
name: "Asha"                     user: "u1"  (ObjectId ref User)
email: "asha@mail.com"           bio: "Full-stack developer"
                                  location: "Kolkata"
\end{verbatim}
\end{diagrambox}

\begin{lstlisting}[language=JavaScript]
// models/Profile.js
const profileSchema = new mongoose.Schema({
  user: { type: mongoose.Schema.Types.ObjectId, ref: "User", required: true, unique: true },
  bio: String,
  location: String,
});
\end{lstlisting}

\section{One-to-Many: User and Tasks}

\begin{diagrambox}[One-to-Many]
\begin{verbatim}
User "u1"
   |
   |-- Task "t1"  { user: "u1", title: "Fix login bug" }
   |-- Task "t2"  { user: "u1", title: "Write tests"    }
   `-- Task "t3"  { user: "u1", title: "Deploy to prod"  }
\end{verbatim}
\end{diagrambox}

This is exactly the \texttt{Task.user} field from Chapter 13: each Task document stores
a single \texttt{ObjectId} pointing back to its owning User.

\section{populate(): Resolving References}

By default, Mongoose stores and returns only the raw \texttt{ObjectId} for a \texttt{ref}
field. \texttt{.populate()} replaces that id with the actual referenced document.

\begin{lstlisting}[language=JavaScript]
// Without populate
const task = await Task.findById(taskId);

// With populate - only pull name and email from the User document
const task = await Task.find().populate("user", "name email");
\end{lstlisting}

\subsection{Before populate() — raw ObjectId}

\begin{lstlisting}[language=JavaScript]
{
  "_id": "665f1a2b3c4d5e6f7a8b9c0d",
  "title": "Fix login bug",
  "status": "todo",
  "user": "664e0a1b2c3d4e5f6a7b8c9d"
}
\end{lstlisting}

\subsection{After populate("user", "name email")}

\begin{lstlisting}[language=JavaScript]
{
  "_id": "665f1a2b3c4d5e6f7a8b9c0d",
  "title": "Fix login bug",
  "status": "todo",
  "user": {
    "_id": "664e0a1b2c3d4e5f6a7b8c9d",
    "name": "Asha Verma",
    "email": "asha@mail.com"
  }
}
\end{lstlisting}

\begin{notebox}[NOTE]
Passing \texttt{"name email"} as the second argument to \texttt{.populate()} is a projection —
it limits which fields of the referenced document are returned, keeping the payload
small and avoiding accidental exposure of fields like \texttt{password}.
\end{notebox}

\section{Many-to-Many: Users and Projects}

A many-to-many relationship needs an array of references on at least one side —
here, a Project keeps an array of member User ids.

\begin{lstlisting}[language=JavaScript]
// models/Project.js
const projectSchema = new mongoose.Schema(
  {
    name: { type: String, required: true },
    members: [{ type: mongoose.Schema.Types.ObjectId, ref: "User" }],
  },
  { timestamps: true }
);
\end{lstlisting}

\begin{diagrambox}[Many-to-Many]
\begin{verbatim}
Project "Website Redesign"          Project "Mobile App"
   members: [u1, u2, u3]                members: [u2, u4]

User u1 ---belongs to---> Website Redesign
User u2 ---belongs to---> Website Redesign, Mobile App
User u3 ---belongs to---> Website Redesign
User u4 ---belongs to---> Mobile App
\end{verbatim}
\end{diagrambox}

\begin{lstlisting}[language=JavaScript]
const project = await Project.findById(projectId).populate("members", "name email");

// project.members is now an array of user objects, not ids:
// members: [
//   { _id: "u1", name: "Asha", email: "asha@mail.com" },
//   { _id: "u2", name: "Ravi", email: "ravi@mail.com" },
//   ...
// ]
\end{lstlisting}

\begin{tipbox}[TIP]
\texttt{.populate()} works the same way whether the reference field holds a single
\texttt{ObjectId} or an array of them — Mongoose detects the schema type automatically.
\end{tipbox}

\begin{warnbox}[WARNING]
Populate runs an extra query behind the scenes. Populating deeply nested or large
arrays on every request can hurt performance — only populate the fields the response
actually needs.
\end{warnbox}
